\section{Discussion}

\subsection{Asymmetry Between Classifier and Language Claims}

The classifier and the language stages of the pipeline are at different
levels of empirical maturity, and the manuscript reflects that
asymmetry. On the classifier side, the methodology is now substantially
more rigorous than in earlier development phases: the active holdout is
physically separated from the training roots, per-class support is
balanced at $n=10$, the augmented training root has been rebuilt from
the compact root with a class-floor balancing pass, and a fresh v4
checkpoint has been retrained on the rebuilt root and evaluated on the
balanced holdout under calibrated two-view TTA. The classifier results
in \S\ref{sec:results-classifier} are therefore reproducible from the
recorded artifacts.

On the language side, the picture is more partial. The runtime
integration of MiniGPT-v2, the RF-DETR visibility constraint, the
optional retrieval module, and the heuristic hallucination retry are
deployed and exercised on a regular basis. The constrained dual-report
supervision policy has been validated by the focused pilot described in
\S\ref{sec:results-dualreport}. However, the deployed report-writer
checkpoint predates the constrained stage-2 redesign and the balanced
holdout carve, which means any explanation-quality measurement on the
current holdout would be confounded by training-time exposure.
Accordingly, we restrict the language claims of this manuscript to (i) a
description of the runtime integration and (ii) a feasibility result for
the constrained supervision policy. We do not assert post-policy
explanation quality.

\subsection{Decoupling Classifier Quality from Deployed Diagnosis Rate}

The threshold gate of Eq.~\ref{eq:threshold-gate} introduces a
deliberate gap between classifier accuracy on the holdout and the
deployed diagnosis-acceptance rate. A correct prediction can be demoted
to \texttt{unknown} when its top-1 probability falls below
$\tau_{\hat{y}}$, and a marginal correct prediction can be subsequently
demoted by the post-prompt-assembly $0.50$ guard. We do not retune the
thresholds in this manuscript because the deployed runtime now logs the
top-1 and top-2 probabilities to a structured decision log, and we
prefer to retune from real decision data rather than from holdout
evaluation alone. The current manuscript therefore reports classifier
accuracy as a property of the model, not as the deployed user-facing
diagnosis rate.

\subsection{Upstream--Downstream Coupling for Ambiguity}

The constrained dual-report policy underscores a more general property
of the system: usefulness of the language-stage uncertainty behavior is
not determined by the language model alone. A naive ``top-2 implies
plausible alternative'' rule would surface numerically close but
semantically implausible competitors (for example, frost as the rank-2
competitor for healthy on a uniformly green plant), and these would
become legitimate-looking ambiguity reports. The policy becomes
defensible only after both a margin condition
(Eq.~\ref{eq:amb-case-1}--\ref{eq:amb-case-2}) and a curated pair
whitelist (Eq.~\ref{eq:allowed-pairs}) are applied. This is a concrete
instance of a more general principle for systems in which a fluent
generator follows a structured predictor: the generator's uncertainty
expression is only as well-formed as the upstream filter that produces
its inputs.

\subsection{Detector Use as Constraint, Not Localization}

The deployed RF-DETR Small detector reaches a validation
$\mathrm{mAP}_{50}$ of 0.772 and $\mathrm{mAP}_{50:95}$ of 0.598. These
numbers are competitive for a six-class plant-part vocabulary at the
training data scale described in \S\ref{sec:results-detector}, but they
are also not the full criterion for the detector's intended runtime
role. The detector is deployed as a binary visibility oracle, not as a
fine-grained localizer: the report writer consumes the set $V(x)$ of
parts for which any detection clears the per-part threshold $\theta_p$,
not the bounding boxes themselves. We tune $\theta_p$ in favor of recall
on parts whose absence is more frequently asserted than their presence
(in particular \emph{root}). A formal ablation of the visibility
constraint against an unconstrained baseline is identified as future
work.

\subsection{Limits of Local Inference}

A persistent design constraint is that the deployed runtime must be
fully local: no frontier-model API is queried at inference time. Any
quality improvement therefore has to survive deployment without
privileged runtime access to a frontier model. In practice, this places
the burden on stage-2 supervision quality rather than on online
distillation or runtime auxiliary scoring. The constrained dual-report
policy is the principal tool we use to translate frontier-model quality
into local-model quality through the offline supervision channel.

\subsection{Threats to Validity}

We note the following threats to validity.

\paragraph{Small per-class support.}
With $n=10$ per class, both Wilson and bootstrap intervals on per-class
top-1 accuracy are wide (Table~\ref{tab:per-class}). Observed
differences between classes should be treated as suggestive rather than
significant; the aggregate bootstrap interval in
Table~\ref{tab:bootstrap-ci} (80.0--95.7\% top-1) is the appropriate
level at which to compare the deployed checkpoint to alternatives.

\paragraph{Single-distribution evaluation.}
The active holdout reflects a single-cultivar, single-region image
distribution. The manuscript does not include cross-cultivar or
cross-region generalization, and we do not claim it. This is a
deliberate scope restriction: external evaluation is identified as
future work alongside the experiments listed in
\S\ref{sec:results-pending}.

\paragraph{Calibration data overlap.}
The temperature parameter $T$ in Eq.~\ref{eq:tta-calibration} is fit on
the internal validation split of the augmented training root, and ECE
is reported on the disjoint physical holdout. However, both splits are
drawn from the same source distribution, so the reported ECE is an
in-distribution calibration estimate rather than a cross-distribution
estimate.

\paragraph{Detector measurement is firing-rate, not precision/recall.}
The RF-DETR mAP metrics in Table~\ref{tab:rfdetr-final} are computed on
the detector's COCO validation split. We additionally report visibility
firing rates of the deployed detector on the disease-holdout
distribution (Table~\ref{tab:rfdetr-holdout}), but those firing rates
do not constitute precision/recall numbers because the disease holdout
has no plant-part ground truth. Per-disease firing rates are consistent
with what the deployment intent of the detector requires (e.g., 100\%
\emph{root} firing on \emph{root rot}; 80\% \emph{soil} firing on
\emph{overwatered}), but a like-for-like detection benchmark on the
disease distribution would require additional plant-part annotation of
the holdout images.

\paragraph{v3-to-v4 retrain not significant at $n=70$.}
The McNemar paired test in Table~\ref{tab:mcnemar} reports a
two-sided exact $p=0.688$ for the comparison of v3 (85.71\%) and v4
(88.57\%) on the shared 70-image holdout. The retrain is therefore
motivated by split-hygiene reasons (\S\ref{sec:data}) rather than by
a statistically significant accuracy improvement at this sample size.

\paragraph{Pilot scope.}
The dual-report pilot is a 20-image feasibility result. It supports
claims about pipeline executability and policy filtering, but does not
constitute a measurement of post-retraining explanation quality. The
pending experiments in \S\ref{sec:results-pending} are required before
stronger end-to-end claims can be supported.

\paragraph{Heuristic hallucination retry.}
The phrase-trigger retry module is not calibrated against a
false-positive budget; we do not characterize its precision or recall
on any test distribution. It is described in \S\ref{sec:system} as part
of the deployed runtime, but no claim about its effect on report
quality is made.

\subsection{Summary}

The deployed AGAI pipeline supports a stronger classifier claim than
language claim, and the manuscript reports them at the level of evidence
they currently support: the classifier and detector results are
calibrated and measured, while the language stage is described and
shown feasible under the constrained supervision policy without yet
being benchmarked end-to-end.
